% A fragment, as a model would paste into a chat: no preamble.
\section*{Solution to Problem 4}

We need to show that $\sum_{k=1}^{n} k = \frac{n(n+1)}{2}$ for all $n \in \mathbb{N}$.

\textbf{Base case.} For $n = 1$: $1 = \frac{1 \cdot 2}{2}$. \checkmark

\textbf{Inductive step.} Assume it holds for $n$. Then
\begin{align*}
\sum_{k=1}^{n+1} k &= \frac{n(n+1)}{2} + (n+1) \\
&= \frac{(n+1)(n+2)}{2}. \qquad \blacksquare
\end{align*}

\begin{center}
\fbox{\textbf{Answer:} $\dfrac{n(n+1)}{2}$}
\end{center}

\paragraph{Remark.} The formula is due to Gauss.\footnote{Allegedly as a schoolboy.}

\begin{itemize}
    \item \textbf{Time complexity:} $\mathcal{O}(1)$
    \item \textbf{Space complexity:} $\Theta(1)$
\end{itemize}

\hrulefill

\noindent\textit{Note:} \small{This uses \underline{strong} induction only implicitly.}\normalsize

\begin{table}[h!]
\centering
\begin{tabular}{|c|c|}
\hline
$n$ & $S_n$ \\ \hline
1 & 1 \\ \hline
2 & 3 \\ \hline
\end{tabular}
\caption{First values}
\end{table}
